\SourceLeaf{99}%
führen wollte, statt seines li~$x = \int \frac{dx}{\log x}$ lieber jene gewählt, da eine
einförmige Function immer ohne Vergleich als classischer und einfacher anzusehen ist als
eine vielförmige, zumal da $\log x$ selbst schon eine vielförmige Function ist. $\mathord{.}\,\mathord{.}\,\mathord{.}\,\mathord{.}\,\mathord{.}\,\mathord{.}\,\mathord{.}\,\mathord{.}\,\mathord{.}$

$\mathord{.}\,\mathord{.}\,\mathord{.}\,\mathord{.}$ Übrigens glaube ich, dass die Ausdehnung der Untersuchungen auf imaginäre
Argumente zu höchst interessanten Resultaten Anlass geben wird. Doch möchte ich aus den
oben angeführten Gründen lieber die Function $\int \frac{e^x-1}{x}\,dx$ als
$\int \frac{dx}{\log x}$ wählen, weil ich vermuthe, dass erstere concinnere Resultate geben
wird. So zum Beispiel möchte ich sehr gern wissen, ob jene Function oder, was dasselbe ist,
die Reihe

\[
x + \tfrac14 xx + \tfrac{1}{18}x^3 + \text{etc.}
\]

\noindent für gewisse endliche Werthe von $x$ von der Form $a+bi$ wohl $0$ werden kann. Mit
Gewissheit kann ich es noch nicht behaupten, obwohl es mir sehr wahrscheinlich ist. Gibt es
solche Werthe (dann gewiss unendlich viele), so werden diess sehr merkwürdige Grössen sein,
und die ganze Reihe wird sich in unendliche Factoren der Form $(1+2ax+\beta xx)$ zerlegen
lassen. ---

$\mathord{.}\,\mathord{.}\,\mathord{.}\,\mathord{.}\,\mathord{.}\,\mathord{.}\,\mathord{.}\,\mathord{.}\,\mathord{.}\,\mathord{.}\,\mathord{.}\,\mathord{.}\,\mathord{.}$

\Needspace{14\baselineskip}
\begin{center}
BEMERKUNG.
\end{center}

\noindent\hspace*{\parindent}Diese principielle Erörterung über die Theorie der Functionen einer complexen Variabelen
ist, wie aus dem Texte hervorgeht, einem Briefe an \textsc{Bessel} vom 18.\ December 1811
entnommen, der bereits in dem »Briefwechsel zwischen \textsc{Gauss} und \textsc{Bessel}«
(Leipzig, 1880) veröffentlicht worden ist. Es ist übrigens wahrscheinlich, dass die von
\textsc{Gauss} hier entwickelte Auffassung ihm bereits lange vor Abfassung des fraglichen
Briefes geläufig war. Eben auf Grund dieser Anschauungen wird \textsc{Gauss} z.\,B.\ die
vermuthlich einer weit früheren Periode angehörenden Untersuchungen in der zweiten Hälfte
des Art.\ 12 vom »Arithmetisch-geometrischen Mittel« (ges.\ Werke III p.\ 378) angestellt
haben. \hfill \textsc{Fricke}.
